\documentclass[14pt,a4paper,oneside]{extbook}
\usepackage[utf8]{vietnam}
\usepackage{amsmath,amssymb,mathrsfs,fancyhdr,enumerate,multirow,makecell}
\usepackage{tkz-euclide,tikz-3dplot,tikz,tkz-tab}
\usetikzlibrary{arrows.meta,calc,intersections,angles,quotes,backgrounds,shapes.geometric,patterns,through,shadows,positioning}
\usepackage{yhmath}
\usepackage[top=0.7cm, bottom=0.7cm, left=1.0cm, right=1.0cm, headheight=16pt, headsep=3mm, footskip=5mm]{geometry}
\usepackage[hidelinks,unicode]{hyperref}
\usepackage{qrcode}
\usepackage[many]{tcolorbox}
\usepackage{fontawesome}

\pagestyle{fancy}
\fancyhf{}
\fancyhead[L]{\footnotesize\bfseries\color{teal!85!black}\faBook~SỔ TAY KIẾN THỨC TOÁN 5}
\fancyhead[R]{\footnotesize\bfseries\color{blue!85!black}Lớp Toán Thầy Lĩnh -- loptoanthaylinh.vn}
\fancyfoot[C]{\footnotesize\bfseries\color{gray!80!black}-- Sổ tay kiến thức: Bài 29. Luyện tập chung hình phẳng --}
\renewcommand{\headrulewidth}{0.4pt}

\begin{document}

% --- BANNER TIÊU ĐỀ BÀI HỌC ---
\noindent
\begin{tcolorbox}[
  enhanced,
  colback=teal!8!white,
  colframe=teal!80!black,
  arc=3mm,
  boxrule=1.1pt,
  width=\linewidth,
  halign=center,
  top=2.2mm, bottom=2.2mm,
  left=3mm, right=3mm,
  drop fuzzy shadow
]
  {\large\bfseries\color{teal!95!black}\faBookmark\ BÀI 29. LUYỆN TẬP CHUNG HÌNH PHẲNG (CHỦ ĐỀ 5)}\\[2pt]
  {\scriptsize\color{black!75}\itshape Hệ thống toàn diện công thức Tam giác, Hình thang, Hình tròn và bài toán diện tích tô màu trên 1 trang A4}
\end{tcolorbox}

\vspace{3.2mm}
% --- KHỐI 1: BẢNG VÀNG CÔNG THỨC CHỦ ĐỀ 5 ---
\noindent
\begin{tcolorbox}[
  enhanced,
  title={\footnotesize\bfseries\color{white}\faTable\ 1. Bảng vàng công thức diện tích các hình phẳng Chủ đề 5},
  colback=blue!2!white,
  colframe=blue!75!black,
  arc=2.5mm, boxrule=0.8pt, top=1.8mm, bottom=1.8mm, left=2mm, right=2mm
]
  \scriptsize
  \centering
  \begin{tabular}{|l|c|c|l|}
    \hline
    \textbf{Tên hình} & \textbf{Chu vi} & \textbf{Diện tích ($S$)} & \textbf{Công thức tính ngược} \\ \hline
    \textbf{Tam giác} & Tổng 3 cạnh & $S = \dfrac{a \times h}{2}$ & $h = \dfrac{2S}{a}$ \quad | \quad $a = \dfrac{2S}{h}$ \\ \hline
    \textbf{Hình thang} & Tổng 4 cạnh & $S = \dfrac{(a + b) \times h}{2}$ & $h = \dfrac{2S}{a + b}$ \quad | \quad $a + b = \dfrac{2S}{h}$ \\ \hline
    \textbf{Hình tròn} & $C = d \times 3,14 = 2 \times r \times 3,14$ & $S = r \times r \times 3,14$ & $d = C : 3,14$ \quad | \quad $r \times r = S : 3,14$ \\ \hline
  \end{tabular}
\end{tcolorbox}

\vspace{3mm}

% --- KHỐI 2: DẠNG BÀI TẬP DIỆN TÍCH PHẦN TÔ MÀU ---
\noindent
\begin{minipage}[t]{0.52\textwidth}
  \begin{tcolorbox}[
    enhanced,
    title={\footnotesize\bfseries\color{white}\faPaintBrush\ 2. Phương pháp tính diện tích phần tô màu},
    colback=teal!2!white,
    colframe=teal!75!black,
    arc=2.5mm, boxrule=0.8pt, top=1.8mm, bottom=1.8mm, left=2.2mm, right=2.2mm
  ]
    \footnotesize
    $\bullet$ \textbf{Nguyên tắc bù trừ:}\\
    \begin{tcolorbox}[colback=white, colframe=teal!75!black, boxrule=0.5pt, arc=1.5mm, halign=center, top=1mm, bottom=1mm]
      {\scriptsize\bfseries\color{teal!95!black}$S_{\text{tô màu}} = S_{\text{hình bao ngoài}} - S_{\text{phần trắng}}$}
    \end{tcolorbox}
    \textit{VD:} Hình vuông cạnh $4\text{ cm}$ chứa một hình tròn nội tiếp bán kính $2\text{ cm}$:\\
    $S_{\text{vuông}} = 4 \times 4 = 16\text{ cm}^2$.\\
    $S_{\text{tròn}} = 2 \times 2 \times 3,14 = 12,56\text{ cm}^2$.\\
    $S_{\text{4 góc tô màu}} = 16 - 12,56 = 3,44\text{ cm}^2$.
  \end{tcolorbox}
\end{minipage}%
\hfill
\begin{minipage}[t]{0.46\textwidth}
  \begin{tcolorbox}[
    enhanced,
    title={\footnotesize\bfseries\color{white}\faSliders\ 3. Tỉ số phần trăm diện tích},
    colback=purple!2!white,
    colframe=purple!70!black,
    arc=2.5mm, boxrule=0.8pt, top=1.8mm, bottom=1.8mm, left=2mm, right=2mm
  ]
    \footnotesize
    Nếu cạnh hình vuông tăng $10\%$:\\
    $\bullet$ Cạnh mới $= 110\% = 1,1 \times a$.\\
    $\bullet$ Diện tích mới: $1,1 \times a \times 1,1 \times a = 1,21 \times a \times a = 121\% \times S$.\\
    $\Rightarrow$ Diện tích tăng thêm: $121\% - 100\% = \mathbf{21\%}$!\\
    (Không phải tăng $10\%$!).
  \end{tcolorbox}
\end{minipage}

\vspace{3mm}

% --- KHỐI 3: SƠ ĐỒ TƯ DUY TỔNG HỢP ---
\noindent
\begin{tcolorbox}[
  enhanced,
  title={\footnotesize\bfseries\color{white}\faSitemap\ Sơ đồ tư duy toàn bộ Chủ đề 5: Hình phẳng},
  colback=yellow!2!white,
  colframe=teal!85!black,
  arc=3mm, boxrule=0.9pt, top=1.8mm, bottom=1.8mm, left=1.5mm, right=1.5mm,
  halign=center
]
  \begin{tikzpicture}[
    scale=0.92,
    every node/.style={transform shape},
    root/.style={rectangle, rounded corners=5pt, fill=teal!85!black, text=white, font=\bfseries\footnotesize, inner sep=5pt, align=center},
    b1/.style={rectangle, rounded corners=4pt, fill=blue!12, draw=blue!75!black, text=blue!90!black, font=\bfseries\scriptsize, inner sep=4.5pt, align=center},
    b2/.style={rectangle, rounded corners=4pt, fill=purple!12, draw=purple!75!black, text=purple!90!black, font=\bfseries\scriptsize, inner sep=4.5pt, align=center},
    b3/.style={rectangle, rounded corners=4pt, fill=green!15, draw=green!65!black, text=green!80!black, font=\bfseries\scriptsize, inner sep=4.5pt, align=center}
  ]
    \node[root] (core) at (0, 0.85) {\faShapes\ HÌNH PHẲNG LỚP 5};
    \node[b1] (n1) at (-5.0, -0.45) {\textbf{TAM GIÁC:}\\[2pt]$S = (a \times h) : 2$};
    \node[b2] (n2) at (0, -0.65) {\textbf{HÌNH TRÒN:}\\[2pt]$C = d \times 3,14$; $S = r \times r \times 3,14$};
    \node[b3] (n3) at (5.0, -0.45) {\textbf{HÌNH THANG:}\\[2pt]$S = (a + b) \times h : 2$};

    \draw[teal!85!black, thick, -{Stealth[length=2mm]}] (core) -| (n1);
    \draw[teal!85!black, thick, -{Stealth[length=2mm]}] (core) -- (n2);
    \draw[teal!85!black, thick, -{Stealth[length=2mm]}] (core) -| (n3);
  \end{tikzpicture}
\end{tcolorbox}

\vspace{3mm}

% --- KHỐI 4: QUY TẮC ĐỒNG BỘ ĐƠN VỊ & BẪY TRẮC NGHIỆM ---
\noindent
\begin{minipage}[t]{0.54\textwidth}
  \begin{tcolorbox}[
    enhanced,
    title={\footnotesize\bfseries\color{white}\faCheck\ Checklist 3 bước trước khi đặt bút tính},
    colback=green!3!white,
    colframe=green!60!black,
    arc=2.5mm, boxrule=0.8pt, top=1.8mm, bottom=1.8mm, left=2.2mm, right=2.2mm
  ]
    \footnotesize
    $\square$ Đã đổi tất cả kích thước về cùng một đơn vị đo chưa?\\
    $\square$ Đã nhớ chia đôi đối với hình tam giác và hình thang chưa?\\
    $\square$ Đề bài cho đường kính hay bán kính hình tròn?
  \end{tcolorbox}
\end{minipage}%
\hfill
\begin{minipage}[t]{0.44\textwidth}
  \begin{tcolorbox}[
    enhanced,
    title={\footnotesize\bfseries\color{white}\faExclamationTriangle\ Mẹo tránh bẫy trắc nghiệm},
    colback=red!2!white,
    colframe=red!70!black,
    arc=2.5mm, boxrule=0.8pt, top=1.8mm, bottom=1.8mm, left=2.2mm, right=2.2mm
  ]
    \footnotesize
    $\bullet$ Đề bài cho bán kính tăng gấp $2$ lần thì diện tích tăng gấp mấy lần? Trả lời: $2 \times 2 = \mathbf{4\text{ lần}}$ (không phải $2$ lần!).\\
    $\bullet$ Chu vi thì chỉ tăng gấp $2$ lần.
  \end{tcolorbox}
\end{minipage}

\end{document}
